% ---
% school: "college"
% level: "8eme"
% section: null
% chapter: "العمليات والحساب على الأعداد الكسرية"
% source: "Axiomaths"
% title: "المقارنة والترتيب"
% country: "Tunisie"
% year: 2026
% professor: "Axiomaths"
% difficulty: "Facile"
% points: null
% tags: ["fractions", "comparaison", "ordre"]
% ---
% Exercice 5/18 de la série 8ème « العمليات والحساب على الأعداد الكسرية » (statement + \axsol + solution). Compiler avec XeLaTeX + axiomath.sty (+ axiomaths-logo.png).
\documentclass[12pt,a4paper]{article}
\usepackage[level=8b,color=15803D]{axiomath}
\axsetup{title={العمليات والحساب على الأعداد الكسرية}, corrige}

\begin{document}

\begin{corr}[title={المقارنة والترتيب}, stars=2, source={Axiomaths}]
\begin{enumerate}
  \item قارن في كل حالة بين العددين :
  \begin{enumerate}
    \item $\dfrac{5}{6}$ و $\dfrac{7}{9}$ \qquad
    \item $-\dfrac{3}{4}$ و $-\dfrac{4}{5}$ \qquad
    \item $\dfrac{-5}{7}$ و $\dfrac{1}{3}$
  \end{enumerate}
  \item رتّب الأعداد التالية ترتيبا تصاعديا : $\ \dfrac{3}{4}\ ,\ -\dfrac{5}{6}\ ,\ \dfrac{7}{12}\ ,\ 0\ ,\ -\dfrac{2}{3}$
  \item قارن بين $\dfrac{99}{100}$ و $\dfrac{100}{101}$ دون استعمال الحاسبة.
\end{enumerate}
\axsol
\begin{prop}[المقارنة]
لمقارنة عددين كسريين نوحّد مقاميهما (الموجبين) ثم نقارن بسطيهما.
كل عدد سالب أصغر من كل عدد موجب. ومن بين عددين سالبين الأكبر هو الأقرب إلى الصفر.
\end{prop}
\begin{enumerate}[itemsep=14pt]
  \item نوحّد المقامين في كل حالة :
  \begin{enumerate}[itemsep=10pt]
    \item المقام المشترك هو $18$ :
    \[\dfrac{5}{6}=\dfrac{15}{18}\ ;\qquad\dfrac{7}{9}=\dfrac{14}{18}\]
    بما أن $15>14$ فإن \ib{$\dfrac{5}{6}>\dfrac{7}{9}$}.
    \item المقام المشترك هو $20$ :
    \[-\dfrac{3}{4}=-\dfrac{15}{20}\ ;\qquad-\dfrac{4}{5}=-\dfrac{16}{20}\]
    بما أن $-15>-16$ فإن \ib{$-\dfrac{3}{4}>-\dfrac{4}{5}$}.
    \item $\dfrac{-5}{7}$ سالب و $\dfrac{1}{3}$ موجب، إذن \ib{$\dfrac{-5}{7}<\dfrac{1}{3}$}.
  \end{enumerate}
  \item المقام المشترك هو $12$ :
  \[\dfrac{3}{4}=\dfrac{9}{12}\ ,\quad -\dfrac{5}{6}=-\dfrac{10}{12}\ ,\quad \dfrac{7}{12}\ ,\quad 0\ ,\quad -\dfrac{2}{3}=-\dfrac{8}{12}\]
  نرتب البسطين : $-10<-8<0<7<9$ ، إذن :
  \[\ib{$-\dfrac{5}{6}<-\dfrac{2}{3}<0<\dfrac{7}{12}<\dfrac{3}{4}$}\]
  \item المقام المشترك هو $100\times101=10\,100$ :
  \[\dfrac{99}{100}=\dfrac{99\times101}{100\times101}=\dfrac{9\,999}{10\,100}\]
  \[\dfrac{100}{101}=\dfrac{100\times100}{101\times100}=\dfrac{10\,000}{10\,100}\]
  بما أن $9\,999<10\,000$ فإن \ib{$\dfrac{99}{100}<\dfrac{100}{101}$}.

  \textbf{طريقة ثانية :} $1-\dfrac{99}{100}=\dfrac{1}{100}$ و $1-\dfrac{100}{101}=\dfrac{1}{101}$ ، و $\dfrac{1}{100}>\dfrac{1}{101}$ ، فالعدد الأول أبعد عن $1$ إذن هو الأصغر.
\end{enumerate}
\end{corr}

\end{document}
